\chapter{Development heuristics}
Rules applied throughout development, and the reason for each.

\begin{description}
  \item[Fidelity first] the default engine is the most accurate available (reSIDfp, ares, Nuked, Ootake). A lighter
        engine remains selectable, for slow devices and comparison.
  \item[Exact values] every value shown is the chip state at the current frame. No smoothing, no averaging: an
        averaged value is a value that never existed.
  \item[Graphics when enough] when a curve says it all (filter response), no text repeats it. When text is needed,
        it sits next to what it describes.
  \item[One view per synthesis] each chip family has its own view, built around what makes its sound: wavetable for
        the PC Engine, operators and algorithms for FM, envelopes and echo for the Super Nintendo, filter and
        modulations for the SID.
  \item[Measure before claiming] every fix is checked by a measurement (compared per-voice outputs, render, test)
        before delivery; what could not be checked is reported as such.
  \item[Render before delivery] every visual change is checked on a render generated by the tests before being built
        into an APK.
  \item[Minimal patches to third-party code] changes to imported libraries are kept small, marked in the code
        (``zxtune extension'', ``standalone'') and follow the style of the modified file, so they can be merged back.
  \item[Neutral defaults] features that change the layout (SID samples lane, AdLib rhythm lanes) are off by default.
  \item[Licenses checked on import] all imported code is checked for compatibility with GPL~v3.
  \item[No regression] the full test suite runs before significant deliveries; the only tolerated failures are
        catalogue tests needing network access.
\end{description}
